% Options for packages loaded elsewhere
% Options for packages loaded elsewhere
\PassOptionsToPackage{unicode}{hyperref}
\PassOptionsToPackage{hyphens}{url}
\PassOptionsToPackage{dvipsnames,svgnames,x11names}{xcolor}
%
\documentclass[
  portuguese,
  11pt,
]{article}
\usepackage{xcolor}
\usepackage{amsmath,amssymb}
\setcounter{secnumdepth}{5}
\usepackage{iftex}
\ifPDFTeX
  \usepackage[T1]{fontenc}
  \usepackage[utf8]{inputenc}
  \usepackage{textcomp} % provide euro and other symbols
\else % if luatex or xetex
  \usepackage{unicode-math} % this also loads fontspec
  \defaultfontfeatures{Scale=MatchLowercase}
  \defaultfontfeatures[\rmfamily]{Ligatures=TeX,Scale=1}
\fi
\usepackage{lmodern}
\ifPDFTeX\else
  % xetex/luatex font selection
\fi
% Use upquote if available, for straight quotes in verbatim environments
\IfFileExists{upquote.sty}{\usepackage{upquote}}{}
\IfFileExists{microtype.sty}{% use microtype if available
  \usepackage[]{microtype}
  \UseMicrotypeSet[protrusion]{basicmath} % disable protrusion for tt fonts
}{}
\makeatletter
\@ifundefined{KOMAClassName}{% if non-KOMA class
  \IfFileExists{parskip.sty}{%
    \usepackage{parskip}
  }{% else
    \setlength{\parindent}{0pt}
    \setlength{\parskip}{6pt plus 2pt minus 1pt}}
}{% if KOMA class
  \KOMAoptions{parskip=half}}
\makeatother
% Make \paragraph and \subparagraph free-standing
\makeatletter
\ifx\paragraph\undefined\else
  \let\oldparagraph\paragraph
  \renewcommand{\paragraph}{
    \@ifstar
      \xxxParagraphStar
      \xxxParagraphNoStar
  }
  \newcommand{\xxxParagraphStar}[1]{\oldparagraph*{#1}\mbox{}}
  \newcommand{\xxxParagraphNoStar}[1]{\oldparagraph{#1}\mbox{}}
\fi
\ifx\subparagraph\undefined\else
  \let\oldsubparagraph\subparagraph
  \renewcommand{\subparagraph}{
    \@ifstar
      \xxxSubParagraphStar
      \xxxSubParagraphNoStar
  }
  \newcommand{\xxxSubParagraphStar}[1]{\oldsubparagraph*{#1}\mbox{}}
  \newcommand{\xxxSubParagraphNoStar}[1]{\oldsubparagraph{#1}\mbox{}}
\fi
\makeatother

\usepackage{color}
\usepackage{fancyvrb}
\newcommand{\VerbBar}{|}
\newcommand{\VERB}{\Verb[commandchars=\\\{\}]}
\DefineVerbatimEnvironment{Highlighting}{Verbatim}{commandchars=\\\{\}}
% Add ',fontsize=\small' for more characters per line
\usepackage{framed}
\definecolor{shadecolor}{RGB}{241,243,245}
\newenvironment{Shaded}{\begin{snugshade}}{\end{snugshade}}
\newcommand{\AlertTok}[1]{\textcolor[rgb]{0.68,0.00,0.00}{#1}}
\newcommand{\AnnotationTok}[1]{\textcolor[rgb]{0.37,0.37,0.37}{#1}}
\newcommand{\AttributeTok}[1]{\textcolor[rgb]{0.40,0.45,0.13}{#1}}
\newcommand{\BaseNTok}[1]{\textcolor[rgb]{0.68,0.00,0.00}{#1}}
\newcommand{\BuiltInTok}[1]{\textcolor[rgb]{0.00,0.23,0.31}{#1}}
\newcommand{\CharTok}[1]{\textcolor[rgb]{0.13,0.47,0.30}{#1}}
\newcommand{\CommentTok}[1]{\textcolor[rgb]{0.37,0.37,0.37}{#1}}
\newcommand{\CommentVarTok}[1]{\textcolor[rgb]{0.37,0.37,0.37}{\textit{#1}}}
\newcommand{\ConstantTok}[1]{\textcolor[rgb]{0.56,0.35,0.01}{#1}}
\newcommand{\ControlFlowTok}[1]{\textcolor[rgb]{0.00,0.23,0.31}{\textbf{#1}}}
\newcommand{\DataTypeTok}[1]{\textcolor[rgb]{0.68,0.00,0.00}{#1}}
\newcommand{\DecValTok}[1]{\textcolor[rgb]{0.68,0.00,0.00}{#1}}
\newcommand{\DocumentationTok}[1]{\textcolor[rgb]{0.37,0.37,0.37}{\textit{#1}}}
\newcommand{\ErrorTok}[1]{\textcolor[rgb]{0.68,0.00,0.00}{#1}}
\newcommand{\ExtensionTok}[1]{\textcolor[rgb]{0.00,0.23,0.31}{#1}}
\newcommand{\FloatTok}[1]{\textcolor[rgb]{0.68,0.00,0.00}{#1}}
\newcommand{\FunctionTok}[1]{\textcolor[rgb]{0.28,0.35,0.67}{#1}}
\newcommand{\ImportTok}[1]{\textcolor[rgb]{0.00,0.46,0.62}{#1}}
\newcommand{\InformationTok}[1]{\textcolor[rgb]{0.37,0.37,0.37}{#1}}
\newcommand{\KeywordTok}[1]{\textcolor[rgb]{0.00,0.23,0.31}{\textbf{#1}}}
\newcommand{\NormalTok}[1]{\textcolor[rgb]{0.00,0.23,0.31}{#1}}
\newcommand{\OperatorTok}[1]{\textcolor[rgb]{0.37,0.37,0.37}{#1}}
\newcommand{\OtherTok}[1]{\textcolor[rgb]{0.00,0.23,0.31}{#1}}
\newcommand{\PreprocessorTok}[1]{\textcolor[rgb]{0.68,0.00,0.00}{#1}}
\newcommand{\RegionMarkerTok}[1]{\textcolor[rgb]{0.00,0.23,0.31}{#1}}
\newcommand{\SpecialCharTok}[1]{\textcolor[rgb]{0.37,0.37,0.37}{#1}}
\newcommand{\SpecialStringTok}[1]{\textcolor[rgb]{0.13,0.47,0.30}{#1}}
\newcommand{\StringTok}[1]{\textcolor[rgb]{0.13,0.47,0.30}{#1}}
\newcommand{\VariableTok}[1]{\textcolor[rgb]{0.07,0.07,0.07}{#1}}
\newcommand{\VerbatimStringTok}[1]{\textcolor[rgb]{0.13,0.47,0.30}{#1}}
\newcommand{\WarningTok}[1]{\textcolor[rgb]{0.37,0.37,0.37}{\textit{#1}}}

\usepackage{longtable,booktabs,array}
\usepackage{calc} % for calculating minipage widths
% Correct order of tables after \paragraph or \subparagraph
\usepackage{etoolbox}
\makeatletter
\patchcmd\longtable{\par}{\if@noskipsec\mbox{}\fi\par}{}{}
\makeatother
% Allow footnotes in longtable head/foot
\IfFileExists{footnotehyper.sty}{\usepackage{footnotehyper}}{\usepackage{footnote}}
\makesavenoteenv{longtable}
\usepackage{graphicx}
\makeatletter
\newsavebox\pandoc@box
\newcommand*\pandocbounded[1]{% scales image to fit in text height/width
  \sbox\pandoc@box{#1}%
  \Gscale@div\@tempa{\textheight}{\dimexpr\ht\pandoc@box+\dp\pandoc@box\relax}%
  \Gscale@div\@tempb{\linewidth}{\wd\pandoc@box}%
  \ifdim\@tempb\p@<\@tempa\p@\let\@tempa\@tempb\fi% select the smaller of both
  \ifdim\@tempa\p@<\p@\scalebox{\@tempa}{\usebox\pandoc@box}%
  \else\usebox{\pandoc@box}%
  \fi%
}
% Set default figure placement to htbp
\def\fps@figure{htbp}
\makeatother



\ifLuaTeX
\usepackage[bidi=basic]{babel}
\else
\usepackage[bidi=default]{babel}
\fi
% get rid of language-specific shorthands (see #6817):
\let\LanguageShortHands\languageshorthands
\def\languageshorthands#1{}


\setlength{\emergencystretch}{3em} % prevent overfull lines

\providecommand{\tightlist}{%
  \setlength{\itemsep}{0pt}\setlength{\parskip}{0pt}}



 


\usepackage{booktabs}
\usepackage{caption}
\usepackage{longtable}
\usepackage{colortbl}
\usepackage{array}
\usepackage{anyfontsize}
\usepackage{multirow}
\usepackage[utf8]{inputenc} % Caracteres Especiais
\usepackage{indentfirst} % Identação
\usepackage[T1]{fontenc} % Especifica a codificação de fonte a ser usada
\usepackage[left=2cm, top=2cm, right=2cm, bottom=2cm]{geometry}
\usepackage{graphicx} % Figuras
\usepackage{latexsym, amssymb, amsmath, amsfonts, amsthm} % Formatação Matemática
\usepackage{xcolor}

\renewcommand{\familydefault}{\sfdefault}

\usepackage{float, caption}
\floatplacement{table}{H}
\floatplacement{figure}{H}
\captionsetup[table]{position=above}
\captionsetup[figure]{position=below}
\makeatletter
\@ifpackageloaded{caption}{}{\usepackage{caption}}
\AtBeginDocument{%
\ifdefined\contentsname
  \renewcommand*\contentsname{Índice}
\else
  \newcommand\contentsname{Índice}
\fi
\ifdefined\listfigurename
  \renewcommand*\listfigurename{Lista de Figuras}
\else
  \newcommand\listfigurename{Lista de Figuras}
\fi
\ifdefined\listtablename
  \renewcommand*\listtablename{Lista de Tabelas}
\else
  \newcommand\listtablename{Lista de Tabelas}
\fi
\ifdefined\figurename
  \renewcommand*\figurename{Figura}
\else
  \newcommand\figurename{Figura}
\fi
\ifdefined\tablename
  \renewcommand*\tablename{Tabela}
\else
  \newcommand\tablename{Tabela}
\fi
}
\@ifpackageloaded{float}{}{\usepackage{float}}
\floatstyle{ruled}
\@ifundefined{c@chapter}{\newfloat{codelisting}{h}{lop}}{\newfloat{codelisting}{h}{lop}[chapter]}
\floatname{codelisting}{Listagem}
\newcommand*\listoflistings{\listof{codelisting}{Lista de Listagens}}
\makeatother
\makeatletter
\makeatother
\makeatletter
\@ifpackageloaded{caption}{}{\usepackage{caption}}
\@ifpackageloaded{subcaption}{}{\usepackage{subcaption}}
\makeatother
\usepackage{bookmark}
\IfFileExists{xurl.sty}{\usepackage{xurl}}{} % add URL line breaks if available
\urlstyle{same}
\hypersetup{
  pdftitle={Análise Multivariada II},
  pdfauthor={Breno Cauã Rodrigues da Silva},
  pdflang={pt},
  colorlinks=true,
  linkcolor={blue},
  filecolor={Maroon},
  citecolor={Blue},
  urlcolor={Blue},
  pdfcreator={LaTeX via pandoc}}


\title{Análise Multivariada II}
\usepackage{etoolbox}
\makeatletter
\providecommand{\subtitle}[1]{% add subtitle to \maketitle
  \apptocmd{\@title}{\par {\large #1 \par}}{}{}
}
\makeatother
\subtitle{LISTA V - ANÁLISE DISCRIMINANTE (AD)}
\author{Breno Cauã Rodrigues da Silva}
\date{}
\begin{document}
\maketitle


\section*{\texorpdfstring{Os dados do arquivo ``\texttt{Jogadores.xls}''
foram coletados por G. R. Bryce e R. M. Barker (Apresentados em Rencher,
2002) como parte de um estudo preliminar sobre uma possível relação
entre o design do capacete de futebol americano e lesões no pescoço.
Seis medidas da cabeça foram feitas em cada sujeito. Havia 30 sujeitos
em cada um dos três grupos: jogadores de futebol americano do ensino
médio (grupo 1), jogadores de futebol americano universitário (grupo 2)
e não jogadores de futebol americano (grupo 3). As seis variáveis
são:}{Os dados do arquivo ``Jogadores.xls'' foram coletados por G. R. Bryce e R. M. Barker (Apresentados em Rencher, 2002) como parte de um estudo preliminar sobre uma possível relação entre o design do capacete de futebol americano e lesões no pescoço. Seis medidas da cabeça foram feitas em cada sujeito. Havia 30 sujeitos em cada um dos três grupos: jogadores de futebol americano do ensino médio (grupo 1), jogadores de futebol americano universitário (grupo 2) e não jogadores de futebol americano (grupo 3). As seis variáveis são:}}\label{os-dados-do-arquivo-jogadores.xls-foram-coletados-por-g.-r.-bryce-e-r.-m.-barker-apresentados-em-rencher-2002-como-parte-de-um-estudo-preliminar-sobre-uma-possuxedvel-relauxe7uxe3o-entre-o-design-do-capacete-de-futebol-americano-e-lesuxf5es-no-pescouxe7o.-seis-medidas-da-cabeuxe7a-foram-feitas-em-cada-sujeito.-havia-30-sujeitos-em-cada-um-dos-truxeas-grupos-jogadores-de-futebol-americano-do-ensino-muxe9dio-grupo-1-jogadores-de-futebol-americano-universituxe1rio-grupo-2-e-nuxe3o-jogadores-de-futebol-americano-grupo-3.-as-seis-variuxe1veis-suxe3o}
\addcontentsline{toc}{section}{Os dados do arquivo
``\texttt{Jogadores.xls}'' foram coletados por G. R. Bryce e R. M.
Barker (Apresentados em Rencher, 2002) como parte de um estudo
preliminar sobre uma possível relação entre o design do capacete de
futebol americano e lesões no pescoço. Seis medidas da cabeça foram
feitas em cada sujeito. Havia 30 sujeitos em cada um dos três grupos:
jogadores de futebol americano do ensino médio (grupo 1), jogadores de
futebol americano universitário (grupo 2) e não jogadores de futebol
americano (grupo 3). As seis variáveis são:}

{
\Large\bf
\begin{itemize}
    \item \texttt{Larg\_cab}: Largura da cabeça na dimensão mais larga;
    \item \texttt{Circ\_cab}: Circunferência da cabeça;
    \item \texttt{Frent\_tras}: Medida da frente para trás na altura dos olhos;
    \item \texttt{Olh\_cab}: Medida do olho ao topo da cabeça;
    \item \texttt{Ore\_cab}: Medida da orelha ao topo da cabeça;
    \item \texttt{Larg\_quei}: Largura da mandíbula.
\end{itemize}
}

\section*{1) Realize a análise exploratória de dados, apresentando a
análise geral e por grupos (apresente tabelas e
gráficos).}\label{realize-a-anuxe1lise-exploratuxf3ria-de-dados-apresentando-a-anuxe1lise-geral-e-por-grupos-apresente-tabelas-e-gruxe1ficos.}
\addcontentsline{toc}{section}{1) Realize a análise exploratória de
dados, apresentando a análise geral e por grupos (apresente tabelas e
gráficos).}

A primeira proposta de análise foi elaborada na
Tabela~\ref{tbl-SUMMARY}, que apresenta as medidas de resumo das
variáveis do conjunto de dados \texttt{Jogadores.xls}. A
Tabela~\ref{tbl-SUMMARY} não leva em consideração a variável de grupo.

\begin{table}

\caption{\label{tbl-SUMMARY}Medidas de Resumo para as Variáveis do
Conjunto de Dados \texttt{Jogadores.xlsx}.}

\centering{

\fontsize{12.0pt}{14.4pt}\selectfont
\begin{tabular*}{\linewidth}{@{\extracolsep{\fill}}lcccccccc}
\toprule
Variável & Média & Desvio Padrão & CV & Mínimo & 1º Quartil & Mediana & 3º Quartil & Máximo \\ 
\midrule\addlinespace[2.5pt]
Circu\_cab & 58,04 & 1,88 & 3,25\% & 54,8 & 56,90 & 57,55 & 59,15 & 63,5 \\ 
Frent\_tras & 19,91 & 0,75 & 3,75\% & 18,5 & 19,42 & 19,90 & 20,40 & 21,8 \\ 
Larg\_cab & 15,40 & 0,67 & 4,32\% & 13,5 & 15,00 & 15,50 & 15,65 & 17,5 \\ 
Larg\_quei & 12,00 & 0,64 & 5,30\% & 10,5 & 11,50 & 12,00 & 12,50 & 13,5 \\ 
olh\_cab & 11,37 & 1,68 & 14,76\% & 7,4 & 10,12 & 11,30 & 12,73 & 15,0 \\ 
ore\_cab & 13,96 & 0,96 & 6,85\% & 11,1 & 13,40 & 13,90 & 14,50 & 16,5 \\ 
\bottomrule
\end{tabular*}

}

\end{table}%

Podemos observar que, em termos de dispersão, a variável
\texttt{olh\_cab} (Medida do olho ao topo da cabeça) possui o maior
Desvio Padrão em proporção à sua Média (14,76\%), indicando uma
variabilidade relativamente alta em comparação com outras medidas como
\texttt{Circu\_cab} (Circunferência da cabeça), que tem uma
variabilidade proporcional menor (3,25\%). Outro ponto é a proximidade
entre a Média e a Mediana para a maioria das variáveis (ex:
\texttt{Frent\_tras}, Média=19,91 e Mediana=19,90) sugere distribuições
razoavelmente simétricas.

Uma outra análise visual univariada pode ser feita a partir da
Figura~\ref{fig-HIST_UNI_VAR}. Onde é apresentado o histograma para cada
variável do \emph{dataset} analisado.

\begin{figure}[H]

\begin{minipage}{0.33\linewidth}

\centering{

\pandocbounded{\includegraphics[keepaspectratio]{RLS_LISTA_V_files/figure-pdf/fig-HIST_UNI_VAR-1.pdf}}

}

\subcaption{\label{fig-HIST_UNI_VAR-1}\texttt{Larg\_cab}}

\end{minipage}%
%
\begin{minipage}{0.33\linewidth}

\centering{

\pandocbounded{\includegraphics[keepaspectratio]{RLS_LISTA_V_files/figure-pdf/fig-HIST_UNI_VAR-2.pdf}}

}

\subcaption{\label{fig-HIST_UNI_VAR-2}\texttt{Circu\_cab}}

\end{minipage}%
%
\begin{minipage}{0.33\linewidth}

\centering{

\pandocbounded{\includegraphics[keepaspectratio]{RLS_LISTA_V_files/figure-pdf/fig-HIST_UNI_VAR-3.pdf}}

}

\subcaption{\label{fig-HIST_UNI_VAR-3}\texttt{Frent\_tras}}

\end{minipage}%
\newline
\begin{minipage}{0.33\linewidth}

\centering{

\pandocbounded{\includegraphics[keepaspectratio]{RLS_LISTA_V_files/figure-pdf/fig-HIST_UNI_VAR-4.pdf}}

}

\subcaption{\label{fig-HIST_UNI_VAR-4}\texttt{olh\_cab}}

\end{minipage}%
%
\begin{minipage}{0.33\linewidth}

\centering{

\pandocbounded{\includegraphics[keepaspectratio]{RLS_LISTA_V_files/figure-pdf/fig-HIST_UNI_VAR-5.pdf}}

}

\subcaption{\label{fig-HIST_UNI_VAR-5}\texttt{ore\_cab}}

\end{minipage}%
%
\begin{minipage}{0.33\linewidth}

\centering{

\pandocbounded{\includegraphics[keepaspectratio]{RLS_LISTA_V_files/figure-pdf/fig-HIST_UNI_VAR-6.pdf}}

}

\subcaption{\label{fig-HIST_UNI_VAR-6}\texttt{Larg\_quei}}

\end{minipage}%

\caption{\label{fig-HIST_UNI_VAR}Histograma das Variáveis Quantitativas
do Conjunto de Dados \texttt{Jogadores.xls}.}

\end{figure}%

Os histogramas oferecem uma primeira visão sobre a distribuição de cada
variável, ainda sem separação por grupo. Os histogramas de \textbf{(a)},
\textbf{(c)} e \textbf{(f)} aparentam ser uma forma ligeiramente
simétrica com uma concentração de dados em torno do centro - lembrando
uma distribuição gaussiana.

Apenas as variáveis de \textbf{(d)} e \textbf{(e)} se distanciam desse
realidade de forma mais grosseira, com histogramas que lembram mais uma
\emph{distribuição bimodal} (com dois picos claros), isso é um forte
indicativo de que os dados, quando vistos de forma agregada, escondem a
existência de pelo menos dois subgrupos com características distintas.
Essa observação justifica plenamente a necessidade de analisar os dados
por grupo. Uma conclusão feita em \textbf{(b)} é que a variável
\texttt{Circu\_cab} mostra uma leve assimetria à direita.

Visando uma análise global juntamente da finalidade deste relatório, foi
elaborada a Tabela~\ref{tbl-SUMMARY_BY_GROUP}, que recalcula as medidas
de resumo por grupo para cada uma das variáveis.

\begin{table}

\caption{\label{tbl-SUMMARY_BY_GROUP}Medidas de Resumo para as Variáveis
do Conjunto de Dados \texttt{Jogadores.xlsx} por Grupo.}

\centering{

\fontsize{12.0pt}{14.4pt}\selectfont
\begin{tabular*}{\linewidth}{@{\extracolsep{\fill}}ccccccccc}
\toprule
Grupo & Média & Desvio Padrão & CV & Mínimo & 1º Quartil & Mediana & 3º Quartil & Máximo \\ 
\midrule\addlinespace[2.5pt]
\multicolumn{9}{l}{Circu\_cab} \\[2.5pt] 
\midrule\addlinespace[2.5pt]
1 & 58,95 & 2,05 & 3,48\% & 55,9 & 57,20 & 58,40 & 60,67 & 63,5 \\ 
2 & 57,39 & 1,71 & 2,98\% & 54,8 & 56,50 & 57,00 & 58,00 & 62,6 \\ 
3 & 57,77 & 1,55 & 2,68\% & 55,0 & 56,75 & 57,65 & 58,65 & 61,7 \\ 
\midrule\addlinespace[2.5pt]
\multicolumn{9}{l}{Frent\_tras} \\[2.5pt] 
\midrule\addlinespace[2.5pt]
1 & 20,11 & 0,84 & 4,19\% & 19,0 & 19,50 & 20,00 & 20,88 & 21,8 \\ 
2 & 19,80 & 0,74 & 3,75\% & 18,5 & 19,30 & 19,90 & 20,20 & 21,5 \\ 
3 & 19,81 & 0,62 & 3,11\% & 18,7 & 19,50 & 19,80 & 20,38 & 20,8 \\ 
\midrule\addlinespace[2.5pt]
\multicolumn{9}{l}{Larg\_cab} \\[2.5pt] 
\midrule\addlinespace[2.5pt]
1 & 15,20 & 0,74 & 4,86\% & 13,5 & 15,00 & 15,50 & 15,50 & 17,5 \\ 
2 & 15,42 & 0,64 & 4,13\% & 14,3 & 15,12 & 15,40 & 15,50 & 17,3 \\ 
3 & 15,58 & 0,58 & 3,71\% & 14,6 & 15,30 & 15,50 & 15,90 & 17,3 \\ 
\midrule\addlinespace[2.5pt]
\multicolumn{9}{l}{Larg\_quei} \\[2.5pt] 
\midrule\addlinespace[2.5pt]
1 & 12,27 & 0,69 & 5,64\% & 10,5 & 12,00 & 12,25 & 12,50 & 13,5 \\ 
2 & 11,94 & 0,61 & 5,14\% & 10,5 & 11,53 & 11,90 & 12,38 & 13,3 \\ 
3 & 11,80 & 0,52 & 4,41\% & 10,8 & 11,50 & 11,75 & 12,10 & 13,3 \\ 
\midrule\addlinespace[2.5pt]
\multicolumn{9}{l}{olh\_cab} \\[2.5pt] 
\midrule\addlinespace[2.5pt]
1 & 13,08 & 1,04 & 7,97\% & 10,0 & 12,50 & 13,00 & 13,88 & 15,0 \\ 
2 & 10,08 & 1,07 & 10,65\% & 8,2 & 9,45 & 9,90 & 10,47 & 12,8 \\ 
3 & 10,95 & 1,21 & 11,03\% & 7,4 & 10,40 & 10,90 & 11,78 & 13,1 \\ 
\midrule\addlinespace[2.5pt]
\multicolumn{9}{l}{ore\_cab} \\[2.5pt] 
\midrule\addlinespace[2.5pt]
1 & 14,73 & 0,94 & 6,41\% & 13,0 & 14,00 & 14,50 & 15,50 & 16,5 \\ 
2 & 13,45 & 0,76 & 5,61\% & 11,1 & 12,85 & 13,45 & 14,05 & 14,5 \\ 
3 & 13,70 & 0,63 & 4,57\% & 12,0 & 13,40 & 13,65 & 13,97 & 15,1 \\ 
\bottomrule
\end{tabular*}

}

\end{table}%

A Tabela~\ref{tbl-SUMMARY_BY_GROUP} aprofunda a análise ao calcular as
mesmas medidas de resumo, mas agora segmentadas para cada um dos três
grupos: jogadores do ensino médio (1), jogadores universitários (2) e
não jogadores (3). As princiapais diferenças observadas:

\begin{itemize}
\item
  As variáveis \textbf{\texttt{olh\_cab}} (Olho ao topo da cabeça),
  \textbf{\texttt{ore\_cab}} (Orelha ao topo da cabeça) e
  \textbf{\texttt{Circu\_cab}} (Circunferência da cabeça) apresentam
  diferenças entre os grupos ao olharmos para média e mediana de cada
  uma delas. Sendo que existe um padrão nessa diferença de grupos, pois,
  o Grupo 1 apresenta estatísticas maiores que o Grupo 3 que, por sua
  vez, apresenta estatísticas maiores que o Grupo 2.
\item
  Em contra-partida, \textbf{\texttt{Frent\_tras}} (Frente para trás),
  \textbf{\texttt{Larg\_cab}} (Largura da cabeça na dimensão mais larga)
  e \textbf{\texttt{Larg\_quei}} (argura da mandíbula) mostraram
  estatísticas muito próximas entre os três grupos, indicando que ela
  provavelmente tem baixo poder para discriminar os grupos.
\end{itemize}

No entanto, é importante ressaltar que as conclusões são baseadas em uma
análise descritiva dos dados. Tendo em vista que as disparidades
existentes entre os grupos não são de alta magnitude em todas as
variáveis, para afirmar que os grupos de fato diferem entre si,
poderíamos aplicar, por exemplo, uma MANOVA.

Por fim, foi desenhada a Figura~\ref{fig-BOXPLOT_BY_GROUP}. Podemos
avaliar a diferênça entre os grupos dentro das variáveis de uma forma
visual por meio de Boxplots.

\begin{figure}[H]

\centering{

\pandocbounded{\includegraphics[keepaspectratio]{RLS_LISTA_V_files/figure-pdf/fig-BOXPLOT_BY_GROUP-1.pdf}}

}

\caption{\label{fig-BOXPLOT_BY_GROUP}Boxplot das Variáveis do Conjunto
de Dados \texttt{Jogadores.xlsx} por Grupo.}

\end{figure}%

Os boxplots da Figura~\ref{fig-BOXPLOT_BY_GROUP} permitem uma comparação
visual direta da distribuição das variáveis entre os grupos, confirmando
e reforçando as observações da análise das medidas de resumo da
Tabela~\ref{tbl-SUMMARY_BY_GROUP}. As variáveis
\textbf{\texttt{olh\_cab}}, \textbf{\texttt{ore\_cab}} e
\textbf{\texttt{Circu\_cab}} demonstram as diferenças mais nítidas: as
``caixas'' do Grupo 1 estão posicionadas visivelmente acima das caixas
dos Grupos 2 e 3, indicando pouca sobreposição. Visualmente,
\textbf{\texttt{olh\_cab}} parece ser a variável com maior poder de
discriminação. Essa clara separação entre os grupos para ambas as
variáveis explica a bimodalidade - das variáveis
\textbf{\texttt{olh\_cab}} e \textbf{\texttt{ore\_cab}} - e assimetria -
da variável \textbf{\texttt{Circu\_cab}} - observada nos histogramas da
Figura~\ref{fig-HIST_UNI_VAR}.

Em contrapartida, as demais variáveis mostram um poder de distinção
menor. \textbf{\texttt{Frent\_tras}}, por exemplo, apresenta boxplots
quase perfeitamente alinhados, confirmando sua baixa capacidade de
diferenciar os grupos. As outras medidas também exibem sobreposição
considerável, apesar de algumas diferenças nas medianas e na presença de
valores discrepantes (outliers) no gráfico.

De modo a fechar a análise exploratória e confirmar tudo o que já foi
dito, foi elaborado a Figura~\ref{fig-PAIRPLOT}.

\begin{figure}[H]

\centering{

\pandocbounded{\includegraphics[keepaspectratio]{RLS_LISTA_V_files/figure-pdf/fig-PAIRPLOT-1.pdf}}

}

\caption{\label{fig-PAIRPLOT}Pairplot do Conjunto de Dados Segmentadas
pelo Grupo de Jogadores.}

\end{figure}%

\begin{center}\rule{0.5\linewidth}{0.5pt}\end{center}

\begin{quote}
\textbf{NOTA:} A título de bônus para a Questão 2, e como complemento à
Análise Exploratória de Dados, será executada uma Análise de Variância
Multivariada (MANOVA) com o objetivo de verificar a existência de
diferenças estatisticamente significativas entre os grupos.
\end{quote}

\begin{center}\rule{0.5\linewidth}{0.5pt}\end{center}

\section*{2) Faça testes de normalidade multivariada e teste de
homogeneidade de
variância.}\label{fauxe7a-testes-de-normalidade-multivariada-e-teste-de-homogeneidade-de-variuxe2ncia.}
\addcontentsline{toc}{section}{2) Faça testes de normalidade
multivariada e teste de homogeneidade de variância.}

A Análise Discriminante apresenta os seguintes pressupostos:

\begin{enumerate}
\def\labelenumi{\arabic{enumi}.}
\tightlist
\item
  Normalidade Multivariada das observações;
\item
  Matrizes de Covariâncias dos grupos deve ser iguais.
\end{enumerate}

Desta forma, para verificar a \emph{Normalidade Multivariada}, foram
aplicados os Testes de Normalidade de Mardia, Henze-Zirklers e de
Royston sob as hipóteses:

\[
\begin{cases} H_{0}: \text{Os dados seguem uma distribuição normal multivariada;} \\ H_{1}: \text{Os dados não seguem uma distribuição normal multivariada.} \end{cases}
\]

\noindent Veja a Tabela~\ref{tbl-MULT_NORM_TEST_MARDIA},
Tabela~\ref{tbl-MULT_NORM_TEST_HZ},
Tabela~\ref{tbl-MULT_NORM_TEST_ROYSTON}.

\begin{longtable}[]{@{}
  >{\centering\arraybackslash}p{(\linewidth - 8\tabcolsep) * \real{0.2237}}
  >{\centering\arraybackslash}p{(\linewidth - 8\tabcolsep) * \real{0.2895}}
  >{\centering\arraybackslash}p{(\linewidth - 8\tabcolsep) * \real{0.1447}}
  >{\centering\arraybackslash}p{(\linewidth - 8\tabcolsep) * \real{0.1579}}
  >{\centering\arraybackslash}p{(\linewidth - 8\tabcolsep) * \real{0.1842}}@{}}

\caption{\label{tbl-MULT_NORM_TEST_MARDIA}Teste de Mardia para
Normalidade Multivariada.}

\tabularnewline

\toprule\noalign{}
\begin{minipage}[b]{\linewidth}\centering
Teste
\end{minipage} & \begin{minipage}[b]{\linewidth}\centering
Estatística de Teste
\end{minipage} & \begin{minipage}[b]{\linewidth}\centering
Valor \(p\)
\end{minipage} & \begin{minipage}[b]{\linewidth}\centering
Método
\end{minipage} & \begin{minipage}[b]{\linewidth}\centering
Normalidade
\end{minipage} \\
\midrule\noalign{}
\endhead
\bottomrule\noalign{}
\endlastfoot
Mardia Skewness & 84,869 & 0,008 & asymptotic & Not normal \\
Mardia Kurtosis & 0,927 & 0,354 & asymptotic & Normal \\

\end{longtable}

\begin{longtable}[]{@{}
  >{\centering\arraybackslash}p{(\linewidth - 8\tabcolsep) * \real{0.2055}}
  >{\centering\arraybackslash}p{(\linewidth - 8\tabcolsep) * \real{0.3014}}
  >{\centering\arraybackslash}p{(\linewidth - 8\tabcolsep) * \real{0.1507}}
  >{\centering\arraybackslash}p{(\linewidth - 8\tabcolsep) * \real{0.1644}}
  >{\centering\arraybackslash}p{(\linewidth - 8\tabcolsep) * \real{0.1781}}@{}}

\caption{\label{tbl-MULT_NORM_TEST_HZ}Teste de Henze-Zirkler para
Normalidade Multivariada.}

\tabularnewline

\toprule\noalign{}
\begin{minipage}[b]{\linewidth}\centering
Teste
\end{minipage} & \begin{minipage}[b]{\linewidth}\centering
Estatística de Teste
\end{minipage} & \begin{minipage}[b]{\linewidth}\centering
Valor \(p\)
\end{minipage} & \begin{minipage}[b]{\linewidth}\centering
Método
\end{minipage} & \begin{minipage}[b]{\linewidth}\centering
Normalidade
\end{minipage} \\
\midrule\noalign{}
\endhead
\bottomrule\noalign{}
\endlastfoot
Henze-Zirkler & 0,941 & 0,228 & asymptotic & Normal \\

\end{longtable}

\begin{longtable}[]{@{}
  >{\centering\arraybackslash}p{(\linewidth - 8\tabcolsep) * \real{0.1324}}
  >{\centering\arraybackslash}p{(\linewidth - 8\tabcolsep) * \real{0.3235}}
  >{\centering\arraybackslash}p{(\linewidth - 8\tabcolsep) * \real{0.1618}}
  >{\centering\arraybackslash}p{(\linewidth - 8\tabcolsep) * \real{0.1765}}
  >{\centering\arraybackslash}p{(\linewidth - 8\tabcolsep) * \real{0.2059}}@{}}

\caption{\label{tbl-MULT_NORM_TEST_ROYSTON}Teste de Royston para
Normalidade Multivariada.}

\tabularnewline

\toprule\noalign{}
\begin{minipage}[b]{\linewidth}\centering
Teste
\end{minipage} & \begin{minipage}[b]{\linewidth}\centering
Estatística de Teste
\end{minipage} & \begin{minipage}[b]{\linewidth}\centering
Valor \(p\)
\end{minipage} & \begin{minipage}[b]{\linewidth}\centering
Método
\end{minipage} & \begin{minipage}[b]{\linewidth}\centering
Normalidade
\end{minipage} \\
\midrule\noalign{}
\endhead
\bottomrule\noalign{}
\endlastfoot
Royston & 30,218 & \textless0.001 & asymptotic & Not normal \\

\end{longtable}

De certa forma houve empate entre os testes apresentados pela
Tabela~\ref{tbl-MULT_NORM_TEST_MARDIA},
Tabela~\ref{tbl-MULT_NORM_TEST_HZ},
Tabela~\ref{tbl-MULT_NORM_TEST_ROYSTON}. Para corrigir isso foi,
primeiramente, realizado o Teste de Shapiro-Wilk para normalidade
univariada. Os resultados estão disponíveis na
Tabela~\ref{tbl-SHAPIRO_UNIVAR}.

\begin{table}

\caption{\label{tbl-SHAPIRO_UNIVAR}Teste de Shapiro-Wilk para
Normalidade Univariada.}

\centering{

\fontsize{12.0pt}{14.4pt}\selectfont
\begin{tabular*}{\linewidth}{@{\extracolsep{\fill}}lccl}
\toprule
Variável & Estatística \(W\) & Valor \(p\) & Normalidade \\ 
\midrule\addlinespace[2.5pt]
Larg\_cab & 0,9464 & 0,0010 & Not normal \\ 
Circu\_cab & 0,9453 & 0,0009 & Not normal \\ 
Frent\_tras & 0,9797 & 0,1723 & Normal \\ 
olh\_cab & 0,9832 & 0,2967 & Normal \\ 
ore\_cab & 0,9727 & 0,0548 & Normal \\ 
Larg\_quei & 0,9857 & 0,4338 & Normal \\ 
\bottomrule
\end{tabular*}

}

\end{table}%

Foi constatada a não normalidade apenas às variáveis \texttt{Larg\_cab}
e \texttt{Circu\_cab}. Com isso, foi aplicada tranformação de Box-Cox
com o parâmetro \(\lambda\) estimado por máxima verossimilhança
(\(\lambda\) ótimo), dada pela seguinte expressão:
\[Y^{(\lambda)} = \dfrac{Y^{\lambda} - 1}{\lambda}.\] Foi usada a
seguinte sintaxe no R:

\begin{Shaded}
\begin{Highlighting}[]
\CommentTok{\# {-}{-}{-}{-}{-}{-}{-}{-}{-}{-}{-}{-}{-}{-}{-}{-}{-}{-}{-}{-}{-}{-}{-}{-}{-}{-}{-}}
\CommentTok{\# [2.2] TRANSFORMAÇÃO BOX{-}COX}
\CommentTok{\# {-}{-}{-}{-}{-}{-}{-}{-}{-}{-}{-}{-}{-}{-}{-}{-}{-}{-}{-}{-}{-}{-}{-}{-}{-}{-}{-}}
\NormalTok{Players.trf }\OtherTok{\textless{}{-}} \DecValTok{1} \SpecialCharTok{*}\NormalTok{ Players}

\NormalTok{m1.for.boxcox }\OtherTok{\textless{}{-}} \FunctionTok{lm}\NormalTok{(}\AttributeTok{formula =}\NormalTok{ Larg\_cab }\SpecialCharTok{\textasciitilde{}} \DecValTok{1}\NormalTok{, }\AttributeTok{data =}\NormalTok{ Players.trf)}
\NormalTok{r1.boxcox }\OtherTok{\textless{}{-}}\NormalTok{ MASS}\SpecialCharTok{::}\FunctionTok{boxcox}\NormalTok{(m1.for.boxcox, }\AttributeTok{lambda =} \FunctionTok{seq}\NormalTok{(}\SpecialCharTok{{-}}\DecValTok{4}\NormalTok{, }\DecValTok{4}\NormalTok{, }\AttributeTok{by =} \FloatTok{0.01}\NormalTok{))}
\NormalTok{lambda1.optim }\OtherTok{\textless{}{-}}\NormalTok{ r1.boxcox}\SpecialCharTok{$}\NormalTok{x[}\FunctionTok{which.max}\NormalTok{(r1.boxcox}\SpecialCharTok{$}\NormalTok{y)]}
\NormalTok{Players.trf}\SpecialCharTok{$}\NormalTok{Larg\_cab }\OtherTok{\textless{}{-}}\NormalTok{ (Players.trf}\SpecialCharTok{$}\NormalTok{Larg\_cab}\SpecialCharTok{\^{}}\NormalTok{lambda1.optim }\SpecialCharTok{{-}} \DecValTok{1}\NormalTok{) }\SpecialCharTok{/}\NormalTok{ lambda1.optim}

\NormalTok{m2.for.boxcox }\OtherTok{\textless{}{-}} \FunctionTok{lm}\NormalTok{(}\AttributeTok{formula =}\NormalTok{ Circu\_cab }\SpecialCharTok{\textasciitilde{}} \DecValTok{1}\NormalTok{, }\AttributeTok{data =}\NormalTok{ Players.trf)}
\NormalTok{r2.boxcox }\OtherTok{\textless{}{-}}\NormalTok{ MASS}\SpecialCharTok{::}\FunctionTok{boxcox}\NormalTok{(m2.for.boxcox, }\AttributeTok{lambda =} \FunctionTok{seq}\NormalTok{(}\SpecialCharTok{{-}}\DecValTok{4}\NormalTok{, }\DecValTok{4}\NormalTok{, }\AttributeTok{by =} \FloatTok{0.01}\NormalTok{))}
\NormalTok{lambda2.optim }\OtherTok{\textless{}{-}}\NormalTok{ r1.boxcox}\SpecialCharTok{$}\NormalTok{x[}\FunctionTok{which.max}\NormalTok{(r2.boxcox}\SpecialCharTok{$}\NormalTok{y)]}
\NormalTok{Players.trf}\SpecialCharTok{$}\NormalTok{Circu\_cab }\OtherTok{\textless{}{-}}\NormalTok{ (Players.trf}\SpecialCharTok{$}\NormalTok{Circu\_cab}\SpecialCharTok{\^{}}\NormalTok{lambda2.optim }\SpecialCharTok{{-}} \DecValTok{1}\NormalTok{) }\SpecialCharTok{/}\NormalTok{ lambda2.optim}
\end{Highlighting}
\end{Shaded}

Desta forma, foi aplicado novamente o Teste de Shapiro-Wilk para
normalidade univariada com resultados dispostos na
Tabela~\ref{tbl-SHAPIRO_UNIVAR2}. E posteriomente os Testes de
Normalidade Multivariada apresentados na
Tabela~\ref{tbl-MULT_NORM_TEST_MARDIA2},
Tabela~\ref{tbl-MULT_NORM_TEST_HZ2},
Tabela~\ref{tbl-MULT_NORM_TEST_ROYSTON2}.

\begin{table}

\caption{\label{tbl-SHAPIRO_UNIVAR2}Teste de Shapiro-Wilk para
Normalidade Univariada.}

\centering{

\fontsize{12.0pt}{14.4pt}\selectfont
\begin{tabular*}{0,7\linewidth}{@{\extracolsep{\fill}}lccl}
\toprule
Variável & Estatística \(W\) & Valor \(p\) & Normalidade \\ 
\midrule\addlinespace[2.5pt]
Larg\_cab & 0,9530 & 0,0026 & Not normal \\ 
Circu\_cab & 0,9746 & 0,0754 & Normal \\ 
Frent\_tras & 0,9797 & 0,1723 & Normal \\ 
olh\_cab & 0,9832 & 0,2967 & Normal \\ 
ore\_cab & 0,9727 & 0,0548 & Normal \\ 
Larg\_quei & 0,9857 & 0,4338 & Normal \\ 
\bottomrule
\end{tabular*}

}

\end{table}%

\begin{longtable}[]{@{}
  >{\centering\arraybackslash}p{(\linewidth - 8\tabcolsep) * \real{0.2267}}
  >{\centering\arraybackslash}p{(\linewidth - 8\tabcolsep) * \real{0.2933}}
  >{\centering\arraybackslash}p{(\linewidth - 8\tabcolsep) * \real{0.1467}}
  >{\centering\arraybackslash}p{(\linewidth - 8\tabcolsep) * \real{0.1600}}
  >{\centering\arraybackslash}p{(\linewidth - 8\tabcolsep) * \real{0.1733}}@{}}

\caption{\label{tbl-MULT_NORM_TEST_MARDIA2}Teste de Mardia para
Normalidade Multivariada para os Dados Transformados.}

\tabularnewline

\toprule\noalign{}
\begin{minipage}[b]{\linewidth}\centering
Teste
\end{minipage} & \begin{minipage}[b]{\linewidth}\centering
Estatística de Teste
\end{minipage} & \begin{minipage}[b]{\linewidth}\centering
Valor \(p\)
\end{minipage} & \begin{minipage}[b]{\linewidth}\centering
Método
\end{minipage} & \begin{minipage}[b]{\linewidth}\centering
Normalidade
\end{minipage} \\
\midrule\noalign{}
\endhead
\bottomrule\noalign{}
\endlastfoot
Mardia Skewness & 73,966 & 0,054 & asymptotic & Normal \\
Mardia Kurtosis & 0,504 & 0,615 & asymptotic & Normal \\

\end{longtable}

\begin{longtable}[]{@{}
  >{\centering\arraybackslash}p{(\linewidth - 8\tabcolsep) * \real{0.2055}}
  >{\centering\arraybackslash}p{(\linewidth - 8\tabcolsep) * \real{0.3014}}
  >{\centering\arraybackslash}p{(\linewidth - 8\tabcolsep) * \real{0.1507}}
  >{\centering\arraybackslash}p{(\linewidth - 8\tabcolsep) * \real{0.1644}}
  >{\centering\arraybackslash}p{(\linewidth - 8\tabcolsep) * \real{0.1781}}@{}}

\caption{\label{tbl-MULT_NORM_TEST_HZ2}Teste de Henze-Zirkler para
Normalidade Multivariada para os Dados Transformados.}

\tabularnewline

\toprule\noalign{}
\begin{minipage}[b]{\linewidth}\centering
Teste
\end{minipage} & \begin{minipage}[b]{\linewidth}\centering
Estatística de Teste
\end{minipage} & \begin{minipage}[b]{\linewidth}\centering
Valor \(p\)
\end{minipage} & \begin{minipage}[b]{\linewidth}\centering
Método
\end{minipage} & \begin{minipage}[b]{\linewidth}\centering
Normalidade
\end{minipage} \\
\midrule\noalign{}
\endhead
\bottomrule\noalign{}
\endlastfoot
Henze-Zirkler & 0,921 & 0,357 & asymptotic & Normal \\

\end{longtable}

\begin{longtable}[]{@{}
  >{\centering\arraybackslash}p{(\linewidth - 8\tabcolsep) * \real{0.1324}}
  >{\centering\arraybackslash}p{(\linewidth - 8\tabcolsep) * \real{0.3235}}
  >{\centering\arraybackslash}p{(\linewidth - 8\tabcolsep) * \real{0.1618}}
  >{\centering\arraybackslash}p{(\linewidth - 8\tabcolsep) * \real{0.1765}}
  >{\centering\arraybackslash}p{(\linewidth - 8\tabcolsep) * \real{0.2059}}@{}}

\caption{\label{tbl-MULT_NORM_TEST_ROYSTON2}Teste de Royston para
Normalidade Multivariada para os Dados Transformados.}

\tabularnewline

\toprule\noalign{}
\begin{minipage}[b]{\linewidth}\centering
Teste
\end{minipage} & \begin{minipage}[b]{\linewidth}\centering
Estatística de Teste
\end{minipage} & \begin{minipage}[b]{\linewidth}\centering
Valor \(p\)
\end{minipage} & \begin{minipage}[b]{\linewidth}\centering
Método
\end{minipage} & \begin{minipage}[b]{\linewidth}\centering
Normalidade
\end{minipage} \\
\midrule\noalign{}
\endhead
\bottomrule\noalign{}
\endlastfoot
Royston & 21,261 & 0,002 & asymptotic & Not normal \\

\end{longtable}

Os resultados mostram que nem após a transformação Box-Cox a variável
\texttt{Larg\_cab} aderiu a distribuição normal. Entretanto, os testes
de normalidade se mostraram mais concordantes entre si (adessão dos
dados a normalidade multivariada), sendo que apenas o Teste de Royston
apresentou resultado divergente. Agora, que os dados seguem uma
distribuição normal multivariada, temos que verificar a
\emph{Homogeneidade/Homocedasticidade das Matrizes de Covariâncias}, foi
aplicado o Teste M de Box sob as hipóteses:

\[
\begin{cases} H_{0}: \text{As matrizes de covariâncias são iguais;} \\ H_{1}: \text{As matrizes de covariâncias não são iguais.} \end{cases}
\]

\noindent Veja a Tabela~\ref{tbl-MBox}.

\begin{longtable}[]{@{}ccc@{}}

\caption{\label{tbl-MBox}Teste M de Box para Homogeneidade das Matrizes
de Covariâncias dos Grupos do Conjunto de Dados.}

\tabularnewline

\toprule\noalign{}
Estatística de Teste & Graus de Liberdade & P-Valor \\
\midrule\noalign{}
\endhead
\bottomrule\noalign{}
\endlastfoot
55,79338 & 42 & 0,0753552 \\

\end{longtable}

Como os pressupostos da Análise de Discriminante e Análise de Variância
Multivariada são quase os mesmos, foi feita uma MANOVA para complemento
da Análise Exploratória de Dados e para verificar se existe realmente
diferença entre os grupos de forma estatisticamente significativa. Veja
a Tabela~\ref{tbl-MANOVA}.

\begin{table}

\caption{\label{tbl-MANOVA}Análisede Variância Multivariada para o Fator
Grupo de Jogadores.}

\centering{

\fontsize{12.0pt}{14.4pt}\selectfont
\begin{tabular*}{\linewidth}{@{\extracolsep{\fill}}lrcccccc}
\toprule
Teste & \(gl\) & Estatística & F \(\approx\) & \(gl_{\text{num}}\) & \(gl_{\text{den}}\) & Valor \(p\) & Significância \\ 
\midrule\addlinespace[2.5pt]
Wilks & 1 & 0,54 & 11,68 & 6 & 83 & 1,91 $\times$ 10\textsuperscript{-9} & *** \\ 
Pillai & 1 & 0,46 & 11,68 & 6 & 83 & 1,91 $\times$ 10\textsuperscript{-9} & *** \\ 
Hotelling-Lawley & 1 & 0,84 & 11,68 & 6 & 83 & 1,91 $\times$ 10\textsuperscript{-9} & *** \\ 
Roy & 1 & 0,84 & 11,68 & 6 & 83 & 1,91 $\times$ 10\textsuperscript{-9} & *** \\ 
\bottomrule
\end{tabular*}
\begin{minipage}{\linewidth}
Significância: Rejeita-se \(H_{0}\) para \(\alpha = 0\) ('\texttt{***}'); \(\alpha = 0.001\) ('\texttt{**}'); \(\alpha = 0.01\) ('\texttt{*}'); \(\alpha = 0.05\) ('\texttt{.}')\\
\end{minipage}

}

\end{table}%

Pode-se observar que todos os testes apontam à mesma direção: a
diferença significativa entre os grupos. Logo, prosseguiremos para a
Análise de Discriminante.

\section*{3) Realize a análise discriminante utilizando todas as
variáveis, apresente as funções discriminantes e a avaliação do modelo
com a matriz de confusão e a acurácia (pelo
menos).}\label{realize-a-anuxe1lise-discriminante-utilizando-todas-as-variuxe1veis-apresente-as-funuxe7uxf5es-discriminantes-e-a-avaliauxe7uxe3o-do-modelo-com-a-matriz-de-confusuxe3o-e-a-acuruxe1cia-pelo-menos.}
\addcontentsline{toc}{section}{3) Realize a análise discriminante
utilizando todas as variáveis, apresente as funções discriminantes e a
avaliação do modelo com a matriz de confusão e a acurácia (pelo menos).}

\begingroup
\setlength\LTleft{0,25\linewidth}
\setlength\LTright{0,25\linewidth}\fontsize{12.0pt}{14.4pt}\selectfont

\begin{longtable}{@{\extracolsep{\fill}}lcc}

\caption{\label{tbl-COEFS_LDA}Coeficientes dos Discriminantes Lineares.}

\tabularnewline

\toprule
Variável & LD1 & LD2 \\ 
\midrule\addlinespace[2.5pt]
Larg\_cab & 117,8461 & 173,3276 \\ 
Frent\_tras & 0,0004 & -0,0125 \\ 
olh\_cab & -0,6443 & 0,5461 \\ 
ore\_cab & -0,5119 & -0,3899 \\ 
Larg\_quei & -0,8301 & -1,5138 \\ 
\bottomrule

\end{longtable}

\endgroup

A partir do modelo apresentado pela Tabela~\ref{tbl-COEFS_LDA}, foi
feita uma previsão para os dados e, desta forma, construída a *Matriz de
Confusão dada pela Tabela~\ref{tbl-MATRIX_CONFUSION}.

\begingroup
\setlength\LTleft{0,325\linewidth}
\setlength\LTright{0,325\linewidth}\fontsize{12.0pt}{14.4pt}\selectfont

\begin{longtable}{@{\extracolsep{\fill}}l|ccc}

\caption{\label{tbl-MATRIX_CONFUSION}Matriz de Confunsão do Modelo de
Discriminação Linear.}

\tabularnewline

\toprule
 & Grupo 1 & Grupo 2 & Grupo 3 \\ 
\midrule\addlinespace[2.5pt]
Grupo 1 & 26 & 1 & 3 \\ 
Grupo 2 & 1 & 20 & 9 \\ 
Grupo 3 & 2 & 8 & 20 \\ 
\bottomrule

\end{longtable}

\endgroup

A partir da Tabela~\ref{tbl-MATRIX_CONFUSION}, foi calculada a acuracia
(percentual de acerto) do modelo descrito pela
Tabela~\ref{tbl-COEFS_LDA}. A acurácia obtida foi cerca de 73,33\%.

\section*{4) Realize o procedimento stepwise e selecione um subgrupo de
variáveis.}\label{realize-o-procedimento-stepwise-e-selecione-um-subgrupo-de-variuxe1veis.}
\addcontentsline{toc}{section}{4) Realize o procedimento stepwise e
selecione um subgrupo de variáveis.}

\section*{5) Realize nova análise discriminante utilizando as variáveis
selecionadas, apresente as funções discriminantes e a avaliação do
modelo com a matriz de confusão e a acurácia (pelo
menos).}\label{realize-nova-anuxe1lise-discriminante-utilizando-as-variuxe1veis-selecionadas-apresente-as-funuxe7uxf5es-discriminantes-e-a-avaliauxe7uxe3o-do-modelo-com-a-matriz-de-confusuxe3o-e-a-acuruxe1cia-pelo-menos.}
\addcontentsline{toc}{section}{5) Realize nova análise discriminante
utilizando as variáveis selecionadas, apresente as funções
discriminantes e a avaliação do modelo com a matriz de confusão e a
acurácia (pelo menos).}

\section*{6) Pesquise uma função discriminante quadrática e realize mais
uma AD com essa função, fazendo a avaliação do modelo (se desejar faça
com todas as variáveis e com as selecionadas no stepwise, ou faça
somente com as
selecionadas).}\label{pesquise-uma-funuxe7uxe3o-discriminante-quadruxe1tica-e-realize-mais-uma-ad-com-essa-funuxe7uxe3o-fazendo-a-avaliauxe7uxe3o-do-modelo-se-desejar-fauxe7a-com-todas-as-variuxe1veis-e-com-as-selecionadas-no-stepwise-ou-fauxe7a-somente-com-as-selecionadas.}
\addcontentsline{toc}{section}{6) Pesquise uma função discriminante
quadrática e realize mais uma AD com essa função, fazendo a avaliação do
modelo (se desejar faça com todas as variáveis e com as selecionadas no
stepwise, ou faça somente com as selecionadas).}

\section*{7) Compare as funções discriminantes em relação a matriz de
confusão e acurácia e escolha a que você considera
melhor.}\label{compare-as-funuxe7uxf5es-discriminantes-em-relauxe7uxe3o-a-matriz-de-confusuxe3o-e-acuruxe1cia-e-escolha-a-que-vocuxea-considera-melhor.}
\addcontentsline{toc}{section}{7) Compare as funções discriminantes em
relação a matriz de confusão e acurácia e escolha a que você considera
melhor.}




\end{document}
